%% ============================================================
%%  ABSTRACT (English)
%% ============================================================
\chapter*{Abstract}
\thispagestyle{plain}
\addcontentsline{toc}{chapter}{Abstract}
\selectlanguage{english}

\noindent
This Master's project presents the logical design of an intelligent AI agent
layer for the SprintUp educational platform. The work delivers an architecture
that is independent of any particular technology stack, so that the target
platform can implement it on its own infrastructure.

The document focuses on two core contributions. The first is an indexing
mechanism for existing activities, which lets an agent check whether a similar
activity has already been produced before generating a new one. The second is
the design of a Model Context Protocol (MCP) server that exposes the
read and write operations on the syllabus as a catalog of typed tools, together
with a catalog of agent roles that consume this server, each restricted by a
tool whitelist that materializes the principle of least privilege.

The document is organized in two chapters. The introduction positions the
project in context. Chapter~1 presents the agent orchestration patterns,
retrieval augmentation, and an introduction to the Model Context Protocol.
Chapter~2 details the design: vector-based indexing of pedagogical activities,
the \gls{mcp} protocol as applied to the system, and the catalog of agents
retained for this first version. A general conclusion closes the document.

\vspace{0.5cm}
\noindent
\textbf{Keywords:} Generative AI, Model Context Protocol, semantic indexing,
conversational agents, EdTech.
\selectlanguage{french}
\newpage
